\section{Experimental Design}

\subsection{Experiment 1 --- VM Failure Threshold (answers \sq{1} and \sq{4})}

\textbf{Setup:} Deploy \dagster{} on the UTM \vm{} with DockerRunLauncher (each job = one
Docker container) and PostgreSQL~16 as backend.

\textbf{Procedure:} Run all six workload levels (L1--L6). Each level is repeated three
times with a 60-second cooldown between runs. Record job success rate, mean execution time,
execution time standard deviation, \textsc{cpu} utilisation (via \texttt{psutil}), memory
utilisation, and \mttr{} at each level.

\textbf{Goal:} Identify the exact concurrency level at which success rate drops below 95\%
and document the shape of the degradation curve leading to that point.

\subsection{Experiment 2 --- Kubernetes Overhead and Crossover Point (answers \sq{2}, \sq{3})}

\textbf{Setup:} Kind cluster with Kubernetes Run Launcher. Matching resource requests to
the \vm{} experiment.

\textbf{Procedure:} Run identical workload levels L1--L6, three repetitions each. Record
the same metrics as Experiment~1, plus pod scheduling latency (job submission to pod
Running state) and container startup time (pod Running to first job instruction).

\textbf{Goal:} Calculate per-level deltas in success rate, execution time, and variance
between \vm{} and Kubernetes. Identify the crossover point using the formal definition.

\subsection{Experiment 3 --- Failure Blast Radius Test (answers \sq{2})}

\textbf{Setup:} Level~L3 (4 concurrent jobs) running on both \vm{} and Kubernetes in
separate test runs.

\textbf{Procedure:} After four concurrent jobs are running, deliberately crash one job
mid-execution --- using \texttt{kubectl delete pod} on Kubernetes and \texttt{docker kill}
on the \vm{}. Measure: number of other running jobs affected, time until surviving jobs
recover, their final success rate. Repeated three times per environment.

\textbf{Goal:} Test whether pod isolation's theoretical blast radius reduction holds in
practice under real concurrent workload conditions.
